\documentclass{beamer}
%\mathbbol

\mode<presentation>




\usetheme{Frankfurt}%
\usecolortheme{seagull}
\logo{\includegraphics[height=.25in]{clarksonGreen}}

\definecolor{garnet}{RGB}{136,0,0}
%\definecolor{clarksonGreen}{RGB}{0,71,28}
\definecolor{clarksonGreen}{RGB}{0,52,21}
\setbeamercolor{palette primary}{fg=clarksonGreen,bg=white}
\setbeamercolor{palette secondary}{fg=clarksonGreen,bg=white}
\setbeamercolor{palette tertiary}{fg=clarksonGreen,bg=white}
\setbeamercolor{palette quaternary}{bg=clarksonGreen,fg=white}
\setbeamercolor{block title}{fg=black,bg=black!15}
\setbeamercolor{block body}{fg=black,bg=black!10}
\setbeamercolor{titlelike}{bg=clarksonGreen,fg=white} % parent=palette quaternary}

\newcommand{\half}{\mbox{$\frac{1}{2}$}}
\newcommand{\deltat}{\mbox{$\triangle t$}}
\newcommand{\deltax}{\mbox{$\triangle x$}}
\newcommand{\deltay}{\mbox{$\triangle y$}}

\newcommand{\deriv}[2]{\frac{d}{d#2}#1}
\newcommand{\derivTwo}[2]{\frac{d^2}{d#2^2}#1}

\newcommand{\lp}{\left(}
\newcommand{\rp}{\right)}


\newcommand{\arrayTwo}[4]{
  \left[
  \begin{array}{rr}
    #1 & #2 \\
    #3 & #4
  \end{array}
  \right]
}

\newcommand{\detTwo}[4]{
  \left|
  \begin{array}{rr}
    #1 & #2 \\
    #3 & #4
  \end{array}
  \right|
}


\newcommand{\vecTwo}[2]{
  \left[
  \begin{array}{r}
    #1 \\  #2
  \end{array}
  \right]
}


\newcommand{\stateTwo}[2]{
  \begin{array}{rr}
    \mbox{\fontsize{6}{6}\selectfont $#1$} \\  \mbox{\fontsize{6}{6}\selectfont $#2$}
  \end{array}
}


\newcommand{\arrayThree}[9]{
  \left[
    \begin{array}{rrr}
      #1 & #2 & #3 \\
      #4 & #5 & #6 \\
      #7 & #8 & #9
    \end{array}
  \right]
}

\newcommand{\detThree}[9]{
  \left|
    \begin{array}{rrr}
      #1 & #2 & #3 \\
      #4 & #5 & #6 \\
      #7 & #8 & #9
    \end{array}
  \right|
}


\newcommand{\startRowOpsTwo}{
  \left[
    \begin{array}{rr|rr}
}

\newcommand{\oneRowOpsTwo}[4] {
      #1 & #2 & #3 & #4 \\
}


\newcommand{\startRowOpsThree}{
  \left[
    \begin{array}{rrr|rrr}
}

\newcommand{\oneRowOpsThree}[6] {
      #1 & #2 & #3 & #4 & #5 & #6 \\
}


\newcommand{\startRowFour}{
  \left[
    \begin{array}{rrrr}
}

\newcommand{\oneRowFour}[4] {
      #1 & #2 & #3 & #4 \\
}

\newcommand{\stopRowOps}{
    \end{array}
  \right]
}


\newcommand{\vecThree}[3]{
  \left[
  \begin{array}{r}
    #1 \\  #2 \\ #3
  \end{array}
  \right]
}


\newcommand{\stateThree}[3]{
  \begin{array}{r}
    \mbox{\fontsize{6}{6}\selectfont $#1$} \\  
    \mbox{\fontsize{6}{6}\selectfont $#2$} \\ 
    \mbox{\fontsize{6}{6}\selectfont $#3$}
  \end{array}
}


\begin{document}

\title{Ordinary Differential Equations}
\subtitle{Math 232 - Week 1, Day 1}

\author{Kelly Black}
\institute{Clarkson University}
\date{7 Oct 2011}

\begin{frame}
  \titlepage
\end{frame}

\begin{frame}
  \frametitle{Outline}
  \tableofcontents[pausesection,hideallsubsections]
\end{frame}


\section{Determinants}


\begin{frame}
  \frametitle{Definition of the Determinant}

  The determinant of a $2\times 2$ matrix is defined to be
  \begin{eqnarray*}
    |A| & = & \detTwo{a}{b}{c}{d} \\
    & = & ad - bc
  \end{eqnarray*}

\end{frame}


\begin{frame}
  \frametitle{Why?}

  Suppose I want to solve the following system of equations:
  \begin{eqnarray*}
    \arrayTwo{a}{b}{c}{d} \vecTwo{x}{y} & = & \vecTwo{\#}{\#} \\
    \startRowOpsTwo
    \oneRowOpsTwo{a}{b}{1}{0}
    \oneRowOpsTwo{c}{d}{0}{1}
    \stopRowOps \\
    \stateTwo{}{-cR_1+aR_2}
    \startRowOpsTwo
    \oneRowOpsTwo{a}{b}{1}{0}
    \oneRowOpsTwo{0}{ad-bc}{-c}{a}
    \stopRowOps
  \end{eqnarray*}

  \uncover<2->
  {
    If $ad-bc$ is zero then there is no inverse. The linear system
    does not have a unique solution.
  }

  \uncover<3->
  {
    In general if the determinant of a matrix is zero its inverse does
    not exist. 
  }

\end{frame}


\section{Higher Dimensions}

\begin{frame}
  \frametitle{Higher Dimensions}

  Definition: Given a \textbf{square} matrix:
  \begin{eqnarray*}
    A & = & 
    \left[
      \begin{array}{rrr|r|rr}
        a_{11} & a_{12} & \cdots & a_{1j} & \cdots & a_{1n} \\
        \vdots &       &        & \vdots &        & \vdots \\ \hline
        a_{i1} & a_{i2} & \cdots & a_{ij} & \cdots & a_{in} \\ \hline
        \vdots &       &        & \vdots &        & \vdots \\
        a_{n1} & a_{n2} & \cdots & a_{nj} & \cdots & a_{nn} \\
      \end{array}
    \right]
  \end{eqnarray*}

  The \textbf{minor}, $M_{ij}$ of a matrix is found by removing row
  $i$ and column $j$.

\end{frame}


\begin{frame}
  \frametitle{The Cofactor of a matrix}

  The \textbf{cofactor}, $c_{ij}$, of a matrix is defined to be 
  \begin{eqnarray*}
    c_{ij} & = & (-1)^{i+j}\left| M_{ij} \right|.
  \end{eqnarray*}

  Note:
  \begin{eqnarray*}
    \left| M_{ij} \right|
  \end{eqnarray*}
  is the determinant of $M_{ij}$.

\end{frame}


\begin{frame}
  \frametitle{Determinants in Higher Dimensions}

  The determinant of a square matrix with more than two rows is
  defined to be
  \begin{eqnarray*}
    \left| A \right| & = & \sum^n_{j=1} a_{ij} c_{ij}.
  \end{eqnarray*}
  This is the determinant by expansion along the $i^{\mathrm{th}}$
  row.

  An equivalent definition:
  \begin{eqnarray*}
    \left| A \right| & = & \sum^n_{i=1} a_{ij} c_{ij}.
  \end{eqnarray*}
  This is the determinant by expansion along the $j^{\mathrm{th}}$
  column.

  

\end{frame}


\section{Examples of Determinants}

\begin{frame}
  \frametitle{Example}

  \begin{eqnarray*}
    \mathrm{det}
    \arrayThree{3}{2}{-1}{4}{6}{2}{1}{3}{1}
    & = & 
    3\detTwo{6}{2}{3}{1} - 2 \detTwo{4}{2}{1}{1} 
    + (-1) \detTwo{4}{6}{1}{3} \\
    & = & 3 (6-6) - 2(4-2) + (-1) (12-6) \\
    & = & -10
  \end{eqnarray*}

\end{frame}


\begin{frame}
  \frametitle{}

  \begin{eqnarray*}
    \mathrm{det}
    \startRowFour
    \oneRowFour{4}{3}{-1}{0}
    \oneRowFour{1}{0}{1}{7}
    \oneRowFour{0}{1}{2}{6}
    \oneRowFour{2}{0}{3}{2}
    \stopRowOps
    & = & 
    -3 \detThree{1}{1}{7}{0}{2}{6}{2}{3}{2}
    + 0 \detThree{4}{-1}{0}{0}{2}{6}{2}{3}{2} \\
    & & 
    - 1 \detThree{4}{-1}{0}{1}{1}{7}{2}{3}{2}
    + 0 \detThree{4}{-1}{0}{1}{1}{7}{0}{2}{6} \\
    & = & 178
  \end{eqnarray*}


\end{frame}


\section{Properties of Determinants}

\begin{frame}
  \frametitle{Properties of Determinants}

  \begin{eqnarray*}
    \mathrm{det}(AB) & = & \mathrm{det}(A) \cdot \mathrm{det}(B) \\
    \mathrm{det}(I_n) & = & 1
  \end{eqnarray*}

  Note:
  \begin{eqnarray*}
    A A^{-1} & = & I_n \\
    \mathrm{det}(A A^{-1}) & = & \mathrm{det}(I_n) \\
    \mathrm{det}(A) \mathrm{det}(A^{-1}) & = & 1 \\
    \mathrm{det}(A^{-1}) & = & \frac{1}{\mathrm{det}(A)}
  \end{eqnarray*}

  Read Cramer's rule pp. 160-162.

\end{frame}


\section{Vector Spaces}

\begin{frame}
  \frametitle{Vector Spaces}

  Definition of a vector space, V:
  \begin{itemize}
  \item Elements in V are called ``vectors.''
  \item If $\vec{x}$, $\vec{y}$ are members of V then
    $\vec{x}+\vec{y}$ is in V.
  \item If $c$ is a real number and $\vec{x}$ is in V then so is
    $c\vec{x}$.
  \item There is a ``zero vector'' where $\vec{x}+\vec{0}=\vec{x}$ for
    every member of V.
  \item For any $\vec{x}$ in V there is another $\vec{y}$ in V where
    $\vec{x}+\vec{y}=\vec{0}$. ($\vec{y}$ is called ``$-\vec{x}$.'')
  \end{itemize}

  See p. 168 for definition and properties.

\end{frame}

\begin{frame}
  \frametitle{Example}

  $\mathbb{R}^2$ is a vector space.
  \begin{eqnarray*}
    \vecTwo{x}{y} & \in & \mathbb{R}^2 \\
    \vecTwo{x}{y} + \vecTwo{u}{v} & = & \vecTwo{x+u}{y+v} \\
    c \vecTwo{x}{y} & = & \vecTwo{cx}{cy} \\
    \vec{0} & = & \vecTwo{0}{0} \\
    \vecTwo{x}{y} + \vecTwo{-x}{-y} & = & \vec{0}
  \end{eqnarray*}

\end{frame}


\begin{frame}
  \frametitle{Example - Continuous Functions}

  The set of continuous functions ($C_0$) is a vector space.
  Suppose $f(t),~g(t)~\in~C_o$:
  \begin{itemize}
  \item $f(t)+g(t)$ is continuous.
  \item $c f(t)$ is continuous if $c$ is a real number.
  \item $h(t)=0$ is the zero vector.
  \item For any function
    \begin{eqnarray*}
      f(t) + (-1)f(t) & = & (1 + (-1))f(t), \\
      & = & 0 f(t), \\
      & = & 0.
    \end{eqnarray*}
  \end{itemize}

\end{frame}


\begin{frame}
  \frametitle{Example - Quadratic Functions}

  The set of quadratic functions is a vector space:
  \begin{eqnarray*}
    h(x) & = & a x^2 + bx + c, \\
    g(x) & = & dx^2 + ex + f,
  \end{eqnarray*}
  where $a$, $b$, $c$, $d$, $e$, and $f$ are real numbers.

  Addition:
  \begin{eqnarray*}
    h(x)+g(x) & = & (a+d) x^2 + (b+e) x + (c+f)
  \end{eqnarray*}
  is a quadratic function.

  Scalar multiplication, if $r$ is a real number:
  \begin{eqnarray*}
    r h(x) & = & ra x^2 + rb x + rc
  \end{eqnarray*}
  is a quadratic.

\end{frame}


\begin{frame}
  \frametitle{Counter Example}

  The set of straight lines through $(0,1)$ is not a vector space:
  \begin{eqnarray*}
    h(x) & = & ax + 1 \\
    g(x) & = & bx + 1
  \end{eqnarray*}

  It fails the addition test:
  \begin{eqnarray*}
    h(x) + g(x) & = & (a+b)x + 2
  \end{eqnarray*}
  is a straight line but does not go through $(0,1)$.

\end{frame}




\end{document}

% LocalWords:  Clarkson pausesection hideallsubsections
